\chapter{Hello, \LaTeX{}!}
\label{chap:hello}
\index{\LaTeX{}!advantages}\index{source-to-PDF workflow}\index{WYSIWYM}

Imagine that you are preparing a report with three equations, a table, and a
figure. You insert a new paragraph near the beginning. The figure moves, a
heading remains at the bottom of a page, and every later figure number must be
checked. The scientific argument has not changed, but the page has demanded
your attention.

\LaTeX{} offers a different division of labour. You describe the document's
content and structure in a plain-text source file. A program interprets those
instructions and builds the typeset document. You still make design decisions,
but you make them as rules and relationships rather than repeated adjustments
to individual pages.

\begin{learningoutcomes}
  \item Explain what \LaTeX{} is and what it is not.
  \item Distinguish a source file from a compiled PDF.
  \item Describe the job of a compiler, document class, and package.
  \item Compare a structured \LaTeX{} workflow with a word-processor workflow.
  \item Identify situations in which \LaTeX{} is useful and situations in which
        another tool may be more practical.
  \item Describe the learning route from a minimal document to a scholarly
        manuscript, thesis, and report.
\end{learningoutcomes}

\section{A document preparation system}

\LaTeX{} is a \term{document preparation system}: a collection of rules and
programs for turning structured instructions into typeset output. It is widely
used for scientific and technical communication, particularly when a document
contains mathematics, references, numbered elements, or many interdependent
parts. The \LaTeX{} Project describes it as a system for preparing documents and
provides the core documentation maintained with the software
\parencite{latexprojectDocumentation}.

\LaTeX{} is not simply a font menu or a single visual editor. The familiar part
of the work is writing sentences. The unfamiliar part is that some plain-text
instructions tell \LaTeX{} what those sentences mean. For example, a heading is
identified as a heading rather than made larger by hand. A figure is given a
caption and a permanent label rather than assigned a typed number.

\begin{classroomnote}
The name is pronounced in several ways. You may hear something close to
``lay-tech'' or ``lah-tech''. Pronunciation does not affect the work. The
capitalisation is written \emph{\LaTeX{}} in ordinary prose.
\end{classroomnote}

\section{Two files with different jobs}

A beginning \LaTeX{} project normally has at least two forms of the document:

\begin{description}[style=nextline,leftmargin=1.6em]
  \item[The source file]
  A plain-text file such as \file{main.tex}. It contains your words and \LaTeX{}
  instructions. You edit this file.

  \item[The output file]
  A typeset file such as \file{main.pdf}. It contains the pages you read,
  print, share, or submit. You inspect this file, but you normally do not use a
  PDF editor to change its content.
\end{description}

The source is a little like a recipe: it states the ingredients and the
operations. The PDF is the prepared result. The analogy has limits---a PDF can
be built repeatedly from the same source without consuming anything---but it
helps us keep the two files' roles separate.

\begin{figure}[htbp]
  \centering
  \begin{tikzpicture}[
  node distance=1.15cm,
  every node/.style={font=\small},
  filebox/.style={
    draw=black,
    line width=0.75pt,
    fill=white,
    minimum width=2.4cm,
    minimum height=1.2cm,
    align=center
  },
  processbox/.style={
    draw=black,
    line width=0.75pt,
    fill=white,
    minimum width=2.4cm,
    minimum height=1.2cm,
    align=center
  },
  flow/.style={-{Stealth[length=2.2mm]},line width=0.75pt,color=black}
]
  \node[filebox] (source) {Source\\\file{main.tex}};
  \node[processbox,right=of source] (compiler) {\LaTeX{}\\compiler};
  \node[filebox,right=of compiler] (pdf) {Output\\\file{main.pdf}};
  \draw[flow] (source) -- node[above]{reads} (compiler);
  \draw[flow] (compiler) -- node[above]{writes} (pdf);
  \draw[flow] (pdf.south) .. controls +(0,-0.9) and +(0,-0.9) ..
    node[midway,below=2pt]{inspect and revise} (source.south);
\end{tikzpicture}

  \caption{The basic \LaTeX{} cycle. We edit the source, compile it, inspect the
  PDF, and revise the source.}
  \label{fig:source-to-pdf}
\end{figure}

The program that reads the source and creates output is a \term{compiler}.
The act of running it is \term{compilation}. As \cref{fig:source-to-pdf}
shows, compilation is not a final ceremony. It is part of a short cycle:
write a small amount, compile, inspect, and improve.

\begin{pausepredict}
Suppose you find a misspelt word in the PDF. Which file should you edit: the
PDF or \file{main.tex}? Predict what must happen after the edit before reading
on.

You edit \file{main.tex}, compile again, and inspect the newly generated PDF.
\end{pausepredict}

\section{A first look at source}

The next chapter will build and explain the smallest useful document. For now,
look at its shape:

\begin{latexcode}{Preview: a minimal \LaTeX{} document}
\documentclass{article}
\begin{document}
Hello, \LaTeX{}!
\end{document}
\end{latexcode}

The lines beginning with a backslash are \term{commands}: instructions for
\LaTeX{}. The word in braces after \cmd{documentclass} selects a broad type of
document. Here, \code{article} supplies general rules suitable for a short
scholarly document. The paired \cmd{begin}\required{document} and
\cmd{end}\required{document} lines mark the region whose contents should
appear on the page. The ordinary sentence between them becomes ordinary
typeset text.

\expectedoutput

After successful compilation, the PDF contains a single page with the words
``Hello, \LaTeX{}!'' near the upper-left area of the text block. It does not show
the commands. They guided the construction of the page.

\begin{tryit}
Without compiling, point to the line that selects the document type, the line
that starts the visible content, the visible sentence, and the line that ends
the visible content. This simple source-reading habit will make later error
messages easier to investigate.
\end{tryit}

\section{Class, packages, and content}

Three layers will recur throughout the book.

\begin{enumerate}
  \item A \term{document class} supplies the broad architecture. The standard
        classes \code{article}, \code{report}, and \code{book} make different
        structural assumptions.
  \item A \term{package} is a specialised toolbox added for a purpose. A
        package might manage units, improve tables, or include graphics.
  \item Your content includes the prose, data, equations, figures, citations,
        and document structure that communicate your work.
\end{enumerate}

We will not load a large toolbox before we need it. The first working document
uses only a class and a body. Later, each package will arrive with a problem it
solves, a smallest example, and an explanation.

\begin{scienceinpractice}
A physicist may use the same \code{article} class for equations, a zoologist
for a species description, and a computer scientist for an algorithm. Their
content differs, but the structural ideas---title, sections, figures,
references, and bibliography---remain recognisable. Discipline-specific tools
can be added without rebuilding the document from nothing.
\end{scienceinpractice}

\section{How this differs from a word processor}

The difference is not that one tool has buttons and the other has none. Many
\LaTeX{} editors provide buttons, menus, and PDF previews. The deeper difference
is the working model.

\begin{table}[htbp]
  \centering
  \caption{Two common approaches to preparing a document. The descriptions
  are tendencies, not absolute rules.}
  \label{tab:workflow-comparison}
  \begin{tabularx}{\textwidth}{@{}
    >{\RaggedRight\arraybackslash\bfseries}p{0.23\textwidth}
    >{\RaggedRight\arraybackslash}X
    >{\RaggedRight\arraybackslash}X@{}}
    \toprule
    Question & Visual word-processor approach & Structured \LaTeX{} approach \\
    \midrule
    What do you edit? & A visual representation of the pages & Plain-text
      source that describes content and structure \\
    How is a heading made? & Often by choosing visual properties or a style &
      By identifying text as a section or subsection \\
    How are numbers updated? & Automatically when features are used
      consistently; manually when typed & From labels and references during
      compilation \\
    Where is broad design stored? & In document styles and direct formatting &
      Mainly in the class, packages, and preamble \\
    What is shared for editing? & Usually one document file & Source files,
      bibliography, figures, and build instructions \\
    \bottomrule
  \end{tabularx}
\end{table}

A well-structured word-processor document can also use styles, automatic
captions, and references. A poorly structured \LaTeX{} file can also contain
fragile manual adjustments. The benefit comes from the method, not from the
file extension alone.

\section{Where \LaTeX{} is especially useful}

\LaTeX{} deserves serious consideration when a document has several of the
following features:

\begin{itemize}
  \item mathematical notation or many numbered equations;
  \item a large bibliography with a defined citation style;
  \item figures, tables, sections, and appendices that refer to one another;
  \item repeated structures across many chapters;
  \item a journal, conference, or institution that supplies a \LaTeX{} class;
  \item source that must remain inspectable, portable, and suitable for
        version control;
  \item output that must be rebuilt reliably after revisions.
\end{itemize}

Those strengths help with articles, theses, technical reports, books, lecture
notes, and documentation. They do not mean that \LaTeX{} is automatically the
best tool for every page.

\section{Realistic limitations}

There is an initial learning cost. A missing character can stop compilation,
and an error message may look forbidding before you know which line matters.
Highly visual free-form layouts can take more work than they do in a design
application. Collaborators who do not use \LaTeX{} may need a different review
workflow. Accessibility remediation and publisher conversion still require
care; attractive PDF pages alone do not guarantee an accessible document.

Some tasks are simply quicker elsewhere. A one-page event poster may suit a
graphics editor. A form that colleagues must edit visually may suit a word
processor. A data graphic is usually created in statistical or illustration
software and then included in \LaTeX{}. Good scholarly practice means choosing a
tool for the work, not defending one tool for every purpose.

\begin{goodpractice}
Before adopting a journal or thesis template, verify its current official
instructions. Do not replace its class, bibliography system, margins, or font
choices merely because a personal setup looks more familiar.
\end{goodpractice}

\section{What happens when compilation goes wrong}

The compiler reads literal instructions. If an instruction is incomplete, it
cannot reliably guess what you intended. It writes a \term{log}: a technical
record of the build, including errors and warnings. An editor may display a
shorter explanation beside the source.

\begin{commonerror}
\textbf{Mistake:} treating the first red message as a judgement on the whole
document.

\textbf{What it means:} the compiler has reported a specific point where its
reading failed or became uncertain. Later messages may be consequences of the
first problem.

\textbf{First response:} save the source, read the first meaningful error,
note the reported line, and inspect that line and the few lines before it. We
will practise this routine throughout the book.
\end{commonerror}

Warnings do not always stop a build. A warning may report a line that is too
wide or a reference that needs another pass. Errors and warnings therefore
need interpretation, not panic or automatic deletion.

\section{The journey ahead}

Our sequence is cumulative. First we make one page. Then we understand the
source, organise prose, write mathematics and quantities, build tables and
figures, connect references, and manage a bibliography. Only after those
pieces are familiar do we control large projects and complete the three major
document types.

The end point is independence: you should be able to decide what the document
needs, build the smallest useful version, inspect it, locate reliable package
documentation, and prepare a clean submission without mysterious copied
preamble fragments.

\section{Code lab: source in, PDF out}

\begin{codelab}
Create a new project containing only the file below. Compile it once, change
the title, add a second paragraph, and compile again. Your useful evidence is
the changed PDF---not whether you memorised the commands.
\end{codelab}

\begin{latexcode}{A complete first experiment}
\documentclass{article}

\title{My first \LaTeX{} experiment}
\author{Your Name}

\begin{document}
\maketitle

This sentence came from a source file.

This is a second paragraph.
\end{document}
\end{latexcode}

Now delete the closing brace after \verb|\title{|, compile, and read the first
error. Restore the brace immediately. You have completed the full working
cycle: edit, compile, inspect, diagnose, repair.

\chapterproject

\begin{miniproject}
Create a folder named \file{measurement-methods-article}. Do not add a complex
template. In a plain-text note, write the following three statements:

\begin{enumerate}
  \item The project will compare Method A and Method B using simulated
        demonstration data.
  \item The values do not represent a real experiment or clinical study.
  \item The eventual document will contain an abstract, sections, equations,
        a table, figures, references, an appendix, and submission checks.
\end{enumerate}

This note defines the scholarly task before we choose formatting. In the next
chapter, it becomes your first compilable source file.
\end{miniproject}

\chaptersummary

\LaTeX{} turns structured plain-text source into typeset output. The source and
PDF have different jobs, and a compiler connects them. A document class
provides broad architecture; packages add specialised tools. \LaTeX{} is strong
when a scholarly document has many structured, numbered, or reusable parts,
but it has a learning cost and is not the best tool for every visual task. Our
working habit will be to edit, compile, inspect, and revise in small steps.

\chaptercheck

\begin{chapterchecklist}
  \item I can identify the source file and the output PDF.
  \item I can explain compilation without using unexplained jargon.
  \item I can state the different roles of a class and a package.
  \item I can give one reason to choose \LaTeX{} and one reason to choose another
        tool.
  \item I have created the planning folder and labelled the future data as
        simulated.
\end{chapterchecklist}

\chapterexercisestitle

\begin{chapterexercises}
  \item In one sentence each, define \emph{source file}, \emph{compiler}, and
        \emph{output file}.
  \item A colleague changes a number directly in a PDF and expects the \LaTeX{}
        source to update. Explain the broken link in this workflow.
  \item Read the preview source in this chapter. Predict which line would be
        visible as ordinary text in the PDF.
  \item Name two document features that make a structured workflow valuable.
  \item Classify each item as class, package, content, or output: a results
        paragraph; \code{book}; a bibliography toolbox; \file{report.pdf}.
  \item A collaborator needs a freely positioned one-page infographic and
        will never edit source. Give a reason not to choose \LaTeX{} automatically.
  \item Challenge: describe a document from your field that you would like to
        build by the end of this book. List its recurring structures.
\end{chapterexercises}
